% ===== 第七段：创新点与研究意义 =====
\subsection{本研究的创新点}

针对上述五个研究问题，本课题提出以下两个具有明确贡献边界的创新点：

\begin{enumerate}
  \item \textbf{创新点一：面向LLM Agent工具调用链的多维可解释风险度量框架。}
  针对G1和G2，本课题提出统一日志Schema和多维风险指标体系，
  将工具敏感度$S_{\text{tool}}$、调用深度$D_{\text{call}}$、权限扩散度$P_{\text{diff}}$、
  不可逆操作比例$I_{\text{irrev}}$、敏感数据暴露度$E_{\text{expo}}$、
  人工确认覆盖率$H_{\text{cov}}$等指标纳入统一加权评分框架，并提供指标贡献分解。
  区别于现有以攻击成功率为核心的评估方案\cite{balunovic2024agentdojo,zhang2024injectagent,liu2024formalizing}，
  本框架从序列层面对工具调用链的组合性风险进行可解释量化，
  使"单次调用合法但序列有害"的安全隐患得以被系统性捕获和度量，
  回应了Chu等\cite{chu2026systematic}对T3/T4层面评估空白的明确呼吁。

  \item \textbf{创新点二：基于行为轨迹图的工具滥用检测与反事实归因方法。}
  针对G3和G4，本课题将Agent执行轨迹建模为异构行为轨迹图$G=(V,E,X)$，
  将过去孤立评估单次调用的检测范式扩展至跨步骤、跨权限域的图结构检测。
  提出面向六类工具滥用模式（过度授权、敏感数据外传、不可逆操作缺少确认、
  参数来源污染、跨Agent委托失控、共享记忆污染）的图规则检测方法，
  结合轻量分类器实现组合风险识别；
  进一步提出基于风险路径定位与反事实消融的可解释归因方法，
  输出Top-K风险路径、关键节点贡献和可审计的归因报告，
  填补现有工作\cite{dewitt2026openchallenges,he2025sentinelagent,wang2026authgraph}
  在跨步骤风险路径识别和可解释责任归因方面的空白。
\end{enumerate}

\subsection{研究意义与价值}

\subsubsection{理论意义}

\begin{enumerate}
  \item 本课题从工具调用链出发，将Agent安全问题由"输入输出文本安全"扩展到"过程级行为安全"。
  课题一提出可解释风险度量体系，能够把工具敏感性、权限边界、调用深度、不可逆性和人工确认等因素纳入统一框架；
  课题二进一步将执行轨迹图化，
  把prompt、plan、tool call、observation和final answer之间的依赖关系形式化表达，
  为工具滥用检测和风险归因提供结构化基础。
  该思路有助于弥补现有研究中"重防御、轻度量""重结果、轻过程""重报警、轻归因"的不足。

  \item 本课题将系统安全领域的溯源图思想\cite{inam2023sok,gao2021threatraptor,cheng2024kairos}迁移至LLM Agent场景，
  提出适应Agent执行语义的异构轨迹图表示，
  为Agent可观测性和执行溯源研究提供新的方法论视角，
  有助于推动Agent安全从"攻防对抗"向"过程审计与责任归因"的范式演进。
\end{enumerate}

\subsubsection{实用价值}

\begin{enumerate}
  \item 从工程落地角度看，本课题不要求深度嵌入复杂Agent框架，而是设计一个可插拔审计模块。
  该模块可以接收Agent执行日志，输出工具调用链风险分数、异常模式、关键风险路径和归因报告。
  对于企业办公Agent、数据分析Agent、代码执行Agent和Agentic IoT原型系统而言，
  该审计模块可以作为部署前评估、运行时监控和事后审计的轻量组件，
  回应了OWASP LLM Top 10\cite{owasp2025llmtop10}对工具调用安全审计工具的明确需求。

  \item 本课题构建的低成本可复现任务集与评估框架，
  可为后续Agent安全研究提供标准化的实验基座，
  降低研究门槛，促进该方向的可复现性和可比较性。
\end{enumerate}

\subsubsection{与国家及行业发展战略的契合度}

\begin{enumerate}
  \item 契合《新一代人工智能发展规划》中关于"人工智能安全"的战略部署，
  为AI系统的安全评估与风险防控提供关键技术支撑。
  \item 响应《网络安全法》和《数据安全法》对关键信息基础设施安全保护的要求，
  为LLM Agent在金融、政务等敏感领域的安全部署提供风险度量与审计方法。
  \item 服务于国家人工智能标准化总体组关于"人工智能可信特征"标准的研究，
  为LLM Agent工具调用链的安全性与可解释性评测提供方法论参考。
\end{enumerate}
